%! Populations, Samples, Parameters, and Statistics — Unit 1, Lesson 1.2 — Slides (Beamer)
% PILOT SHAPE (2026-09-06): modeled on AP Statistics lesson 1.4 minus its AP Practice page.
% GRADUAL RELEASE deck: title → targets → hook (unresolved) → I DO divider → two frames per
% notes section → WE DO (a live surface, un-answered) → spoken debrief → close & START THE
% HOMEWORK, which is the individual practice. No group activity, no practice launch, no exit
% ticket. There is no spoiler rule: the notes frames ARE the direct instruction.
\documentclass[aspectratio=169, 11pt]{beamer}
\usepackage{saar-beamer}   % provides \ceruleanheader{} and \sectionlabel[color]{}

\newcommand{\redink}[1]{{\color{redacc}\bfseries #1}}

\begin{document}

% TITLE SLIDE (CourseName is not defined in beamer, so the name is literal)
{
\setbeamercolor{background canvas}{bg=cerulean}
\begin{frame}[plain]
  \vspace{2.8cm}\hspace{0.55cm}%
  \begin{minipage}{0.9\textwidth}
    {\color{goldacc}\bfseries\small LESSON 1.2}\\[0.5cm]
    {\color{white}\bfseries\LARGE Populations, Samples, Parameters, and Statistics}\\[1.2cm]
    {\color{frostmid}\small Statistical Analysis \& Algebraic Reasoning~$\cdot$~Unit 1}
  \end{minipage}
\end{frame}
}

% ── LEARNING TARGETS ─────────────────────────────────────────────────────────
\begin{frame}[t, plain]
  \ceruleanheader{Today's targets}
  \sectionlabel{Lesson 1.2}
  By the end of today, I can\ldots
  \begin{itemize}\setlength{\itemsep}{5pt}
    \item name the \textbf{population} and the \textbf{sample} in a study, and give the sample size.
    \item tell a \textbf{parameter} from a \textbf{statistic} by asking \emph{who the number
          describes}.
    \item name the \textbf{constraint} that keeps a study from being a \textbf{census}.
    \item explain why two honest samples disagree --- \textbf{sampling variability} --- and say
          what a sample does and does not let a report claim.
  \end{itemize}
  \begin{block}{How today works}
    Warm-up $\to$ \textbf{notes together} $\to$ \textbf{one survey we work as a class} $\to$
    debrief $\to$ \textbf{you start the homework here, alone}.
  \end{block}
\end{frame}

% ── WARM-UP ──────────────────────────────────────────────────────────────────
\begin{frame}[t, plain]
  \ceruleanheader{Warm-Up --- 5 minutes, silent}
  \sectionlabel{back to the Lakeside Farmers Market}
  \vspace{0.3cm}
  \begin{center}
    \begin{minipage}{0.86\textwidth}
      About \textbf{800} shoppers come on a Saturday. The manager asked \textbf{50} of them;
      \textbf{26} drove.
      \begin{enumerate}\setlength{\itemsep}{7pt}
        \item What percent of the shoppers she asked drove?
        \item Predict how many of the 800 drove.
        \item Two groups are hiding in that paragraph. How many in each --- and which number
              did she actually \emph{count}?
      \end{enumerate}
    \end{minipage}
  \end{center}
  \vspace{0.4cm}
  \begin{center}
    {\color{cerulean}\bfseries Last night's numbers. Keep the page --- the notes open with them.}
  \end{center}
\end{frame}

% ── HOOK — leave it UNRESOLVED; the second half of section 2 settles it ──────
{
\setbeamercolor{background canvas}{bg=cerulean}
\begin{frame}[plain]
  \vspace{1.3cm}
  \begin{center}
    {\color{frostmid}\bfseries\small THREE VOLUNTEERS. THE SAME SATURDAY MORNING. 50 SHOPPERS EACH.}\\[0.8cm]
    {\color{white}\bfseries\Large Ana got 52\%. \quad Ben got 48\%. \quad Cleo got 58\%.}\\[0.9cm]
    {\color{goldacc}\bfseries\large Who made the mistake?}\\[0.5cm]
    {\color{frostmid}\small Vote: Ana, Ben, Cleo, or nobody. We settle this at the end of
    section 2.}
  \end{center}
\end{frame}
}

% ── I DO — direct instruction begins ─────────────────────────────────────────
{
\setbeamercolor{background canvas}{bg=cerulean}
\begin{frame}[plain]
  \vspace{1.8cm}\hspace{0.8cm}%
  \begin{minipage}{0.86\textwidth}
    {\color{goldacc}\bfseries\small GUIDED NOTES \ \textbullet\ \ I DO}\\[0.7cm]
    {\color{white}\bfseries\Large Open to your Guided Notes page. Fill each blank as we
    name it --- not before.}\\[0.9cm]
    {\color{frostmid}\small Five terms go in the vocabulary box today. Write each one the
    moment we use it.}
  \end{minipage}
\end{frame}
}

% ── 1a. TWO GROUPS ───────────────────────────────────────────────────────────
\begin{frame}[t, plain]
  \ceruleanheader{Two groups}
  \sectionlabel{notes section 1 $\cdot$ population, sample, sample size}
  \vspace{4pt}
  \begin{columns}[t]
    \begin{column}{0.47\textwidth}
      \begin{block}{Population}
        The entire group the question is about --- everyone the conclusion should describe.
      \end{block}
      \vspace{4pt}
      {\small Lakeside: all \textbf{800} Saturday shoppers.}
    \end{column}
    \begin{column}{0.47\textwidth}
      \begin{block}{Sample}
        The part of the population you actually collected data from.
      \end{block}
      \vspace{4pt}
      {\small Lakeside: the \textbf{50} she asked. Sample size $n = 50$.}
    \end{column}
  \end{columns}
  \vspace{12pt}
  \begin{center}
    {\large The population is decided by the \redink{question} --- not by who was handy.}\\[6pt]
    {\small\color{slate} Ask everyone in the population and you have run a \textbf{census}.
    Almost nothing in this course is one.}
  \end{center}
\end{frame}

% ── 1b. PARAMETER AND STATISTIC ──────────────────────────────────────────────
\begin{frame}[t, plain]
  \ceruleanheader{Two kinds of numbers}
  \sectionlabel{notes section 1 $\cdot$ parameter and statistic}
  \vspace{4pt}
  \begin{itemize}\setlength{\itemsep}{6pt}
    \item A number that describes the \textbf{p}opulation is a \textbf{\color{cerulean}parameter}.
          Usually unknown --- and there is only one.
    \item A number that describes the \textbf{s}ample is a \textbf{\color{cerulean}statistic}.
          Known --- and it changes from sample to sample.
    \item To sort a number, ask one question: \redink{who does it describe?}
  \end{itemize}
  \vspace{8pt}
  \begin{block}{Vote first: ``District records show 61\% of Riverbend's 900 students ride a bus.''}
    Parameter or statistic? --- It is a percent. Somebody computed it. And it describes
    \textbf{all 900}. So it is a \redink{parameter}. What the number looks like never decides.
  \end{block}
\end{frame}

% ── 2a. CONSTRAINTS ──────────────────────────────────────────────────────────
\begin{frame}[t, plain]
  \ceruleanheader{Why not just ask all 800?}
  \sectionlabel{notes section 2 $\cdot$ constraints}
  \vspace{2pt}
  \begin{center}
    \small
    \renewcommand{\arraystretch}{1.3}
    \begin{tabular}{@{} l l @{}}
    \hline
    \textbf{Constraint} & \textbf{What it means} \\
    \hline
    Time        & Reaching everyone would take longer than you have. \\
    Cost        & Reaching everyone would cost more than the budget. \\
    Access      & Some individuals cannot be found, reached, or made to answer. \\
    Destruction & Measuring the individual ruins it. \\
    Change      & The population keeps changing while you collect. \\
    \hline
    \end{tabular}
  \end{center}
  \vspace{8pt}
  \begin{center}
    {\color{cerulean}\bfseries A constraint is not an excuse. It is why sampling exists ---
    and part of describing a study honestly.}
  \end{center}
\end{frame}

% ── 2b. THE CRUX ─────────────────────────────────────────────────────────────
\begin{frame}[t, plain]
  \ceruleanheader{The statistic moves. The parameter does not.}
  \sectionlabel[redacc]{notes section 2 $\cdot$ the whole point of today}
  \vspace{4pt}
  \begin{center}
    \small
    \renewcommand{\arraystretch}{1.3}
    \begin{tabular}{@{} l c l @{}}
    \hline
    \textbf{Volunteer} & \textbf{Drove (of 50)} & \textbf{Statistic} \\
    \hline
    Ana  & 26 & $26/50 = 52\%$ \\
    Ben  & 24 & $24/50 = 48\%$ \\
    Cleo & 29 & $29/50 = 58\%$ \\
    \hline
    \end{tabular}
  \end{center}
  \vspace{6pt}
  \begin{itemize}\setlength{\itemsep}{4pt}
    \item Three samples. Three statistics. \textbf{Vote again --- who made the mistake?}
    \item How many true percents are there for all 800 shoppers? \redink{One.} Did it move? \redink{No.}
  \end{itemize}
  \begin{block}{Sampling variability (PS.DC.1e)}
    \textbf{A different answer is not a wrong answer.} A statistic changes from sample to
    sample even though the population did not. Nobody miscounted. \redink{It is what samples do.}
  \end{block}
\end{frame}

% ── WE DO — a LIVE surface: task and data only, no answers ───────────────────
\begin{frame}[t, plain]
  \ceruleanheader{Guided Practice: The Riverbend Job Survey}
  \sectionlabel[goldacc]{We do $\cdot$ you hold the pen}
  The principal wants to know what fraction of Riverbend's \textbf{900} students hold a
  part-time job. She asks \textbf{50}; \textbf{20} have one.
  \begin{itemize}\setlength{\itemsep}{5pt}
    \item[(a)] Population, sample, sample size.
    \item[(b)] The percent, and about how many of the 900 --- then \emph{label} both numbers.
    \item[(c)] The assistant principal asks a \emph{different} 50; \textbf{23} have a job.
          Did somebody make a mistake? Which survey found the true percent?
    \item[(d)] One constraint that kept her from asking all 900.
    \item[(e)] Her report sentence --- then her draft: ``\emph{Exactly} 360 students have a job.''
  \end{itemize}
  \begin{block}{The question behind the question}
    If you picked a survey in (c), you are still at the hook vote.
  \end{block}
\end{frame}

% ── GUIDED PRACTICE (e) — shown AFTER the class has worked it ─────────────────
\begin{frame}[t, plain]
  \ceruleanheader{What a sample lets you say --- and what it does not}
  \sectionlabel[goldacc]{Guided Practice (e) $\cdot$ estimate, not count}
  \vspace{4pt}
  \begin{block}{The principal's sentence}
    ``Of the \textbf{50} students we asked, \textbf{40\%} have a part-time job. If our sample
    represents the school, that is about \textbf{360} of Riverbend's \textbf{900} students.''
  \end{block}
  \vspace{8pt}
  \begin{columns}[t]
    \begin{column}{0.47\textwidth}
      \begin{block}{``Exactly 360 students have a job.''}
        \redink{Not allowed.} 360 is an estimate of the parameter --- nobody counted it.
      \end{block}
    \end{column}
    \begin{column}{0.47\textwidth}
      \begin{block}{``40\% of Riverbend families have a working teenager.''}
        \redink{Not allowed.} Families were never the population.
      \end{block}
    \end{column}
  \end{columns}
\end{frame}

% ── DEBRIEF ──────────────────────────────────────────────────────────────────
\begin{frame}[t, plain]
  \ceruleanheader{Debrief: three statistics, one parameter}
  \sectionlabel{10 minutes $\cdot$ pens down, papers up --- we say these aloud}
  \vspace{4pt}
  \begin{enumerate}\setlength{\itemsep}{5pt}
    \item Ana, Ben, Cleo --- say the sentence. Then: who does the \textbf{61\%} describe?
    \item ``Exactly 360 students have a job.'' Which word is wrong, and why?
  \end{enumerate}
  \vspace{6pt}
  \begin{block}{Cold check --- hands down, thirty seconds}
    A school of 1{,}200 students. The yearbook staff asks 60 and finds 45\% have a job; the
    newspaper asks a different 60 and finds 50\%. A teacher says one of them made an error.
    \textbf{Is the teacher right? What is the true percent?}
  \end{block}
\end{frame}

% ── YOU DO — the homework is the individual practice, started in class ───────
\begin{frame}[t, plain]
  \ceruleanheader{You do: start the homework}
  \sectionlabel[goldacc]{Last 10 minutes $\cdot$ on your own $\cdot$ finish it at home}
  \begin{itemize}\setlength{\itemsep}{5pt}
    \item \textbf{Harbor Point Community Pool} --- 1{,}500 members, 75 asked on a Tuesday evening.
    \item Name the groups; the frequency table; scale 36\% to 1{,}500 \emph{as an estimate}.
    \item Label four numbers --- two of them are percents, and they are not the same kind.
    \item A second survey gets 40\% instead of 36\%. Explain why nobody counted wrong.
  \end{itemize}
  \begin{block}{Silent and alone --- I am circulating. Packet pages, scored out of 12, due at the start of the first class after two study halls.}
    Start with \textbf{1, 4, and 6}. Item 6(b) is the one I am reading over your shoulder ---
    if you blamed a staff member, flag me before you leave.
  \end{block}
\end{frame}

% ── CLOSE ────────────────────────────────────────────────────────────────────
{
\setbeamercolor{background canvas}{bg=cerulean}
\begin{frame}[plain]
  \vspace{1.5cm}\hspace{0.8cm}%
  \begin{minipage}{0.86\textwidth}
    {\color{goldacc}\bfseries\small BEFORE YOU GO}\\[0.7cm]
    {\color{white}\bfseries\Large A statistic describes the sample you have. A parameter
    describes the population you want --- and there is only one. Two honest samples
    disagreeing is not a mistake.}\\[0.9cm]
    {\color{frostmid}\small Next: Lesson 1.3 --- three volunteers each grabbed 50 shoppers.
    \emph{Which} 50? How you choose them decides whether the statistic is worth anything.}
  \end{minipage}
\end{frame}
}

\end{document}
